%2multibyte Version: 5.50.0.2953 CodePage: 1252

\documentclass[notes=show]{beamer}
%%%%%%%%%%%%%%%%%%%%%%%%%%%%%%%%%%%%%%%%%%%%%%%%%%%%%%%%%%%%%%%%%%%%%%%%%%%%%%%%%%%%%%%%%%%%%%%%%%%%%%%%%%%%%%%%%%%%%%%%%%%%%%%%%%%%%%%%%%%%%%%%%%%%%%%%%%%%%%%%%%%%%%%%%%%%%%%%%%%%%%%%%%%%%%%%%%%%%%%%%%%%%%%%%%%%%%%%%%%%%%%%%%%%%%%%%%%%%%%%%%%%%%%%%%%%
\usepackage{amsfonts}
\usepackage{graphicx}
\usepackage{amssymb}
\usepackage{amsmath}
\usepackage{mathpazo}
\usepackage{hyperref}
\usepackage{multimedia}

\setcounter{MaxMatrixCols}{10}
%TCIDATA{OutputFilter=LATEX.DLL}
%TCIDATA{Version=5.50.0.2953}
%TCIDATA{Codepage=1252}
%TCIDATA{<META NAME="SaveForMode" CONTENT="2">}
%TCIDATA{BibliographyScheme=Manual}
%TCIDATA{Created=Wednesday, May 01, 2013 14:54:39}
%TCIDATA{LastRevised=Monday, November 28, 2022 10:22:33}
%TCIDATA{<META NAME="GraphicsSave" CONTENT="32">}
%TCIDATA{<META NAME="DocumentShell" CONTENT="Other Documents\SW\Slides - Beamer">}
%TCIDATA{Language=American English}
%TCIDATA{CSTFile=beamer.cst}
%TCIDATA{PageSetup=72,72,72,83,0}
%TCIDATA{Counters=arabic,1}
%TCIDATA{AllPages=
%H=36
%F=36
%}


\newenvironment{stepenumerate}{\begin{enumerate}[<+->]}{\end{enumerate}}
\newenvironment{stepitemize}{\begin{itemize}[<+->]}{\end{itemize} }
\newenvironment{stepenumeratewithalert}{\begin{enumerate}[<+-| alert@+>]}{\end{enumerate}}
\newenvironment{stepitemizewithalert}{\begin{itemize}[<+-| alert@+>]}{\end{itemize} }
\input{tcilatex}
\usetheme{JuanLesPins}
\usecolortheme{whale}

\begin{document}

\title{Econometrics: Lectures 8\\
Maximum Likelihood Estimation Some Results}
\author{Katerina Petrova \ \ \ \ \ \ \ \ \ \ \ \ \ \ \ \ \ \ \ \ \ \ \ \ \ \
\ \ \ \ \ \ \ \ \ \ \ \ \ \ \ \ \ \ \ \ \ \ \ \ \ \ \ \ \ \ \ \ \ \ \ \ \ \
\ \ \ \ \ \ \ \ \ \ \ \ \ \ \ \ \ \ \ \ \ \ \ }
\date{}
\maketitle

\section{The Linear Gaussian Regression Model}

%TCIMACRO{\TeXButton{BeginFrame}{\begin{frame}}}%
%BeginExpansion
\begin{frame}%
%EndExpansion

\QTR{frametitle}{The Method of Maximum Likelihood}

\QTR{framesubtitle}{Gaussian Linear Model}

\begin{itemize}
\item Ideal Conditions for the Gaussian Linear Model%
\begin{equation*}
y=X\beta +\varepsilon ,\;\;\;\;X\in \mathbb{R}^{n\times k}\text{, }\beta \in 
\mathbb{R}^{k}\text{, }n>k
\end{equation*}

\begin{enumerate}
\item $X$ stochastic matrix of regressors

\item $X$ is full rank: $r\left( X\right) =k$

\item $u|X\sim \mathcal{N}\left( 0,\sigma ^{2}I_{n}\right) $ (independence,
homoskedasticity, normality).
\end{enumerate}

\item Under conditions 1-3, $y|X\sim \mathcal{N}\left( X\beta ,\sigma
^{2}I_{n}\right) $ implying the following likelihood function for a sample
based on $y$:%
\begin{eqnarray*}
L\left( \beta ,\sigma ^{2}\right) &=&f\left( y|X\right) \\
&=&\left( 2\pi \sigma ^{2}\right) ^{-n/2}\exp \left\{ -\frac{1}{2\sigma ^{2}}%
\left( y-X\beta \right) ^{\prime }\left( y-X\beta \right) \right\} .
\end{eqnarray*}
\end{itemize}

%TCIMACRO{\TeXButton{EndFrame}{\end{frame}}}%
%BeginExpansion
\end{frame}%
%EndExpansion

%TCIMACRO{\TeXButton{BeginFrame}{\begin{frame}}}%
%BeginExpansion
\begin{frame}%
%EndExpansion

\QTR{frametitle}{The Method of Maximum Likelihood}

\QTR{framesubtitle}{Gaussian Linear Model}

Therefore, maximising $L\left( \beta ,\sigma ^{2}\right) $ w.r.t. $\beta $
is equivalent to minimising $\left( y-X\beta \right) ^{\prime }\left(
y-X\beta \right) $ w.r.t. $\beta $, i.e. doing OLS. Therefore, for the
regression model under conditions 1-3, \textit{the maximum likelihood
estimator and the OLS estimator of }$\beta $\textit{\ are identical}. The
MLE of $\sigma ^{2}$ can be found as usual as the maximiser of $\log L\left(
\beta ,\sigma ^{2}\right) $ w.r.t. $\sigma ^{2}$:%
\begin{equation*}
\hat{\sigma}_{n}^{2}=\frac{1}{n}\hat{\varepsilon}^{\prime }\hat{\varepsilon}%
,\;\;\text{where\ }\hat{\varepsilon}:=y-X\hat{\beta}\;\;\text{and\ \ }\hat{%
\beta}_{n}=\left( X^{\prime }X\right) ^{-1}X^{\prime }y.
\end{equation*}

%TCIMACRO{\TeXButton{EndFrame}{\end{frame}}}%
%BeginExpansion
\end{frame}%
%EndExpansion

%TCIMACRO{\TeXButton{BeginFrame}{\begin{frame}}}%
%BeginExpansion
\begin{frame}%
%EndExpansion

\QTR{frametitle}{The Method of Maximum Likelihood}

\QTR{framesubtitle}{Gaussian Linear Model}

\begin{eqnarray*}
\frac{\partial ^{2}l(\beta ,\sigma ^{2})}{\partial \beta \partial \beta
^{\prime }} &=&-\frac{1}{\sigma ^{2}}X^{\prime }X \\
\frac{\partial ^{2}l(\beta ,\sigma ^{2})}{\left( \partial \sigma ^{2}\right)
^{2}} &=&\frac{n}{2\sigma ^{4}}-\frac{1}{\sigma ^{6}}(y-X\beta )^{\prime
}(y-X\beta ) \\
\frac{\partial ^{2}l(\beta ,\sigma ^{2})}{\partial \beta \partial \sigma ^{2}%
} &=&\left( \frac{\partial ^{2}l(\beta ,\sigma ^{2})}{\partial \sigma
^{2}\partial \beta ^{\prime }}\right) ^{\prime }=-\frac{1}{\sigma ^{4}}%
(X^{\prime }y-X^{\prime }X\beta )
\end{eqnarray*}

%TCIMACRO{\TeXButton{EndFrame}{\end{frame}}}%
%BeginExpansion
\end{frame}%
%EndExpansion

%TCIMACRO{\TeXButton{BeginFrame}{\begin{frame}}}%
%BeginExpansion
\begin{frame}%
%EndExpansion

\QTR{frametitle}{The Method of Maximum Likelihood}

\QTR{framesubtitle}{Gaussian Linear Model}

The Fisher information matrix is therefore given by%
\begin{gather*}
\mathcal{I}(\beta ,\sigma ^{2})=\left[ 
\begin{array}{cc}
\frac{1}{\sigma ^{2}}\mathbb{E}\left( X^{\prime }X\right) & \frac{1}{\sigma
^{4}}\mathbb{E}(X^{\prime }\varepsilon ) \\ 
\frac{1}{\sigma ^{4}}\mathbb{E}(\varepsilon ^{\prime }X) & -\frac{n}{2\sigma
^{4}}+\frac{1}{\sigma ^{6}}\mathbb{E}(y-X\beta )^{\prime }(y-X\beta )%
\end{array}%
\right] \\
=\left[ 
\begin{array}{cc}
\frac{1}{\sigma ^{2}}\mathbb{E}\left( X^{\prime }X\right) & 0 \\ 
0 & \frac{n}{2\sigma ^{4}}%
\end{array}%
\right]
\end{gather*}%
because $\mathbb{E}(y-X\beta )^{\prime }(y-X\beta )=\mathbb{E}%
\sum_{i=1}^{n}\varepsilon _{i}^{2}=\sum_{i=1}^{n}\mathbb{E[}\varepsilon
_{1}^{2}]=n\sigma ^{2.}$

%TCIMACRO{\TeXButton{EndFrame}{\end{frame}}}%
%BeginExpansion
\end{frame}%
%EndExpansion

%TCIMACRO{\TeXButton{BeginFrame}{\begin{frame}}}%
%BeginExpansion
\begin{frame}%
%EndExpansion

\QTR{framesubtitle}{Gaussian Linear Model}

\begin{itemize}
\item Note that because $\varepsilon |X\sim \mathcal{N}\left( 0,\sigma
^{2}I_{n}\right) ,$ $\hat{\beta}_{n}^{MLE}|X=\left[ \beta +\left( X^{\prime
}X\right) ^{-1}X^{\prime }\varepsilon \right] |X\sim \mathcal{N}\left( \beta
,\sigma ^{2}\left( X^{\prime }X\right) ^{-1}\right) $ even for small $n$

\item This follows since linear functions of normals are also normal

\item To see this, $\hat{\beta}_{n}^{MLE}|X-\beta =\left( X^{\prime
}X\right) ^{-1}X^{\prime }\varepsilon |X\sim \mathcal{N}\left( 0,\left(
X^{\prime }X\right) ^{-1}X^{\prime }\sigma ^{2}X\left( X^{\prime }X\right)
^{-1}\right) $ which under homoskedasticity reduces to $\mathcal{N}\left(
0,\sigma ^{2}\left( X^{\prime }X\right) ^{-1}\right) .$

\item Hence, $\hat{\beta}_{n}^{MLE}|X\sim \mathcal{N}\left( \beta ,\sigma
^{2}\left( X^{\prime }X\right) ^{-1}\right) $ for any sample size

\item In addition, note that for elements of $\hat{\beta}_{j}^{MLE}|X\sim 
\mathcal{N}\left( \beta _{j},\sigma ^{2}\left[ \left( X^{\prime }X\right)
^{-1}\right] _{jj}\right) ,$ so $z=\frac{\widehat{\beta }_{j}-\beta _{j}^{0}%
}{\sqrt{\sigma ^{2}\left[ \left( X^{\prime }X\right) ^{-1}\right] _{jj}}}%
\sim \mathcal{N}(0,1)$
\end{itemize}

%TCIMACRO{\TeXButton{EndFrame}{\end{frame}}}%
%BeginExpansion
\end{frame}%
%EndExpansion

%TCIMACRO{\TeXButton{BeginFrame}{\begin{frame}}}%
%BeginExpansion
\begin{frame}%
%EndExpansion

\QTR{framesubtitle}{Gaussian Linear Model}

\begin{itemize}
\item But $\sigma ^{2}$ is typically not known

\item We have that $\hat{\varepsilon}|X=M_{X}\varepsilon |X\sim \mathcal{N}%
\left( 0,\sigma ^{2}M_{X}\right) $

\item Hence $\widehat{\sigma }_{j}^{2}=\varepsilon ^{\prime
}M_{X}\varepsilon $ turns out to be $\chi ^{2}(n-k)$, where $n-k$ comes from
the rank of $M_{X}$

\item Then $\frac{\widehat{\beta }_{j}-\beta _{j}^{0}}{SE(\widehat{\beta }%
_{j})}$ is not longer Normal, as it is a ratio of two random variables: $%
\mathcal{N}$ and $\chi ^{2}(n-k)$

\item A ratio of $\mathcal{N}(0,1)$ over $\chi ^{2}(n-k)$\ follow a
t-distribution $t(n-k);$ so, we need to use the t critical values instead

\item When $n\rightarrow \infty ,$ the t-distribution is approximately
Normal, so for large enough samples (for $n>30$ the differences are
immaterial) we could use $\mathcal{N}(0,1)$

\item Also note that if $\varepsilon \nsim \mathcal{N},$ we have shown that $%
\widehat{\beta }_{j}$ is asymptotically normal, and $\widehat{\sigma }%
_{j}^{2}$ converges to a point; so the ratio above is asymptotically normal;
for large samples we can still use $\mathcal{N}(0,1)$, even though $%
\varepsilon \nsim \mathcal{N}.$
\end{itemize}

%TCIMACRO{\TeXButton{EndFrame}{\end{frame}}}%
%BeginExpansion
\end{frame}%
%EndExpansion

%TCIMACRO{\TeXButton{BeginFrame}{\begin{frame}}}%
%BeginExpansion
\begin{frame}%
%EndExpansion

\QTR{frametitle}{Testing Linear Restrictions-the F test}

\begin{itemize}
\item What if we want to test a set of linear restrictions

\item Recall that the Wald statistic is $W_{n}=\frac{\left( R\hat{\beta}%
-c\right) ^{\prime }\left[ R\left( X^{\prime }X\right) ^{-1}R^{\prime }%
\right] ^{-1}\left( R\hat{\beta}-c\right) }{\widehat{\sigma }^{2}}$
\end{itemize}

\begin{theorem}
\textit{Consider the linear Gaussian regression model. When }$H_{0}$\textit{%
\ is true,}

$\ \ \ \ \ \ \ \ \ \ \ \ \ \ \ \ \ \ \ \ \ \ \ \ \ \ \ \ \ \frac{1}{r}%
W_{n}\sim F\left( r,n-k\right) $

\textit{where }$F\left( r,n-k\right) $\textit{\ denotes the }$F$\textit{\
distribution with parameters }$r$\textit{\ and }$n-k$\textit{\ and }$r$%
\textit{\ is the rank of }$R$\textit{.}
\end{theorem}

\begin{itemize}
\item This yields critical region for the Wald test based on $W_{n}$: reject 
$H_{0}$ when $W_{n}$ is larger than the F critical value.

\item Unsurprisingly, this test is also known as the \textquotedblleft $F$
test\textquotedblright

\item Whenever $\varepsilon \nsim \mathcal{N}$, we can rely on asymptotic
results to get the distribution of $W_{n}:$ under $H_{0},$ $W_{n}\rightarrow
\chi ^{2}\left( r\right) $ since $\widehat{\sigma }^{2}$ converges to a point%
$\;$as\ $n\rightarrow \infty .$
\end{itemize}

%TCIMACRO{\TeXButton{EndFrame}{\end{frame}}}%
%BeginExpansion
\end{frame}%
%EndExpansion

%TCIMACRO{\TeXButton{BeginFrame}{\begin{frame}}}%
%BeginExpansion
\begin{frame}%
%EndExpansion

\QTR{frametitle}{The T test as a special case of the F test}

\begin{itemize}
\item Also, the t-test is just a special case of the F-test, when we test
one restriction at a time.

\item The (Wald) test with one restriction is%
\begin{equation}
W_{n}=\frac{\left( X_{j}^{\prime }M_{X_{-j}}X_{j}\right) (\hat{\beta}%
_{j}-\beta _{j}^{0})^{2}}{\widehat{s}^{2}}\sim F\left( 1,n-k\right)
\label{W11}
\end{equation}

\item We use the fact that $t_{p}^{2}=F\left( 1,p\right) .$

\item Taking the square root of $W_{n}$ yields an alternative test:%
\begin{equation}
T_{n}:=\frac{\left( X_{j}^{\prime }M_{X_{-j}}X_{j}\right) ^{1/2}(\hat{\beta}%
_{j}-\beta _{j}^{0})}{\widehat{s}}\sim t_{n-k}  \label{t}
\end{equation}%
where $t_{n-k}$ denotes the Student's t-distribution with $n-k$ degrees of
freedom.
\end{itemize}

%TCIMACRO{\TeXButton{EndFrame}{\end{frame}}}%
%BeginExpansion
\end{frame}%
%EndExpansion

\section{Score and Information Matrix}

%TCIMACRO{\TeXButton{BeginFrame}{\begin{frame}}}%
%BeginExpansion
\begin{frame}%
%EndExpansion

\QTR{frametitle}{Score Vector and Information Matrix}

\begin{itemize}
\item Consider a random vector $X=\left( X_{1},...,X_{n}\right) ^{\prime }$
with density $f\left( X,\theta _{0}\right) $ depending on a parameter $%
\theta _{0}\in \Theta \subseteq \mathbb{R}^{k}$.

Let $\theta $ be an arbitrary element of $\Theta $, and let $L\left( \theta
\right) =f\left( X,\theta \right) $ and $l\left( \theta \right) =\log
L\left( \theta \right) $ denote the likelihood function and the
log-likelihood function.
\end{itemize}

\begin{definition}
The $k\times 1$ random vector 
\begin{equation*}
s\left( X,\theta \right) =\partial l/\partial \theta 
\end{equation*}%
is called the \textbf{score vector}.
\end{definition}

\begin{itemize}
\item When $l\left( \theta \right) $ is differentiable, the MLE $\hat{\theta}%
_{n}$ satisfies 
\begin{equation}
s\left( X,\hat{\theta}_{n}\right) =0.  \label{2.12}
\end{equation}
\end{itemize}

%TCIMACRO{\TeXButton{EndFrame}{\end{frame}}}%
%BeginExpansion
\end{frame}%
%EndExpansion

%TCIMACRO{\TeXButton{BeginFrame}{\begin{frame}}}%
%BeginExpansion
\begin{frame}%
%EndExpansion

\QTR{frametitle}{Score Vector and Information Matrix}

First, we will determine the mean and covariance matrix of the score vector $%
s=s\left( X,\theta \right) $.

In what follows, we assume that the following three regularity conditions
are satisfied:

\begin{description}
\item[(RC1)] The support of $f$, $S=\left \{ X\in \mathbb{R}^{n}:f\left(
X,\theta \right) >0\right \} $, is independent of $\theta $.

\item[(RC2)] $f\left( X,\theta \right) $ is twice continuously
differentiable with respect to $\theta .$

\item[(RC3)] The matrix%
\begin{equation}
J\left( \theta \right) :=\mathbb{E}\left( \frac{\partial ^{2}\log f\left(
X,\theta \right) }{\partial \theta \partial \theta ^{\prime }}\right) 
\label{J}
\end{equation}%
is finite $\left( \left\Vert J\left( \theta \right) \right\Vert <\infty
\right) $ and negative definite at $\theta =\theta _{0}$.
\end{description}

%TCIMACRO{\TeXButton{EndFrame}{\end{frame}}}%
%BeginExpansion
\end{frame}%
%EndExpansion

%TCIMACRO{\TeXButton{BeginFrame}{\begin{frame}}}%
%BeginExpansion
\begin{frame}%
%EndExpansion

\QTR{frametitle}{Score Vector and Information Matrix}

\begin{itemize}
\item Condition (RC1) is satisfied for most densities $f\left( \cdot ,\theta
\right) $, e.g., if $X\sim Exp\left( \theta \right) $, then 
\begin{equation*}
f\left( X,\theta \right) =\left\{ 
\begin{array}{ll}
\theta e^{-\theta X}, & X>0 \\ 
0, & X\leq 0.%
\end{array}%
\right.
\end{equation*}%
and $S=\left( 0,\infty \right) $ does not depend on $\theta $.

\item The same holds for all distributions that we have encountered apart
from the uniform distribution on $\left[ 0,\theta \right] $:%
\begin{equation*}
f\left( X,\theta \right) =\left\{ 
\begin{array}{ll}
1/\theta , & X\in \left[ 0,\theta \right] \\ 
0, & X\notin \left[ 0,\theta \right]%
\end{array}%
\right.
\end{equation*}%
so in this case $S=\left[ 0,\theta \right] $ is dependent on $\theta $. This
is why the uniform distribution on $\left[ 0,\theta \right] $ is a source of
counterexamples for likelihood theory.
\end{itemize}

%TCIMACRO{\TeXButton{EndFrame}{\end{frame}}}%
%BeginExpansion
\end{frame}%
%EndExpansion

%TCIMACRO{\TeXButton{BeginFrame}{\begin{frame}}}%
%BeginExpansion
\begin{frame}%
%EndExpansion

\QTR{frametitle}{Score Vector and Information Matrix}

\begin{itemize}
\item The importance of (RC1) lies in that it permits differentiation under
the integral sign, i.e., ensures that%
\begin{equation}
\frac{\partial }{\partial \theta }\int_{S}f\left( X,\theta \right)
dX=\int_{S}\frac{\partial f\left( X,\theta \right) }{\partial \theta }dX,
\label{2.13}
\end{equation}

\item This follows from the so-called Leibnitz rule [see Spiegel, p. 163,
eq. (17)].
\end{itemize}

%TCIMACRO{\TeXButton{EndFrame}{\end{frame}}}%
%BeginExpansion
\end{frame}%
%EndExpansion

%TCIMACRO{\TeXButton{BeginFrame}{\begin{frame}}}%
%BeginExpansion
\begin{frame}%
%EndExpansion

\QTR{frametitle}{Score Vector and Information Matrix}

\begin{theorem}
Under the regularly conditions (RC1)-(RC3),%
\begin{eqnarray}
\mathbb{E}\left[ s\left( X,\theta _{0}\right) \right] &=&0,\   \label{2.14}
\\
\mathbb{E}\left[ s\left( X,\theta _{0}\right) s\left( X,\theta _{0}\right)
^{\prime }\right] &=&-J\left( \theta _{0}\right)
\end{eqnarray}%
where $J\left( \theta \right) $ is the matrix in \emph{(\ref{J})}.
\end{theorem}

%TCIMACRO{\TeXButton{EndFrame}{\end{frame}}}%
%BeginExpansion
\end{frame}%
%EndExpansion

%TCIMACRO{\TeXButton{BeginFrame}{\begin{frame}}}%
%BeginExpansion
\begin{frame}%
%EndExpansion

\QTR{frametitle}{Score Vector and Information Matrix}

Recall that the \textit{true} density of $X$ is $f\left( X,\theta
_{0}\right) $, so that,%
\begin{equation}
\mathbb{E}\left[ g\left( X,\theta \right) \right] =\int_{S}g\left( X,\theta
\right) f\left( X,\theta _{0}\right) dX  \label{moment}
\end{equation}%
for any measurable function $g\left( X,\theta \right) $. Also, $f\left(
X,\theta \right) $ is a density for all $\theta \in \Theta $:%
\begin{equation}
\int_{S}f\left( X,\theta \right) dX=1,\;\;\;\theta \in \Theta .  \label{2.15}
\end{equation}%
Note that $\int_{S}f\left( X,\theta \right) dX$ means%
\begin{equation*}
\underset{\left\{ X\in 
%TCIMACRO{\U{211d} }%
%BeginExpansion
\mathbb{R}
%EndExpansion
^{n}:f\left( X_{1},...,X_{n},\theta \right) >0\right\} }{\int ...\int }%
f\left( X_{1},...,X_{n},\theta \right) dX_{1}....dX_{n}.
\end{equation*}

%TCIMACRO{\TeXButton{EndFrame}{\end{frame}}}%
%BeginExpansion
\end{frame}%
%EndExpansion

%TCIMACRO{\TeXButton{BeginFrame}{\begin{frame}}}%
%BeginExpansion
\begin{frame}%
%EndExpansion

\QTR{frametitle}{Score Vector and Information Matrix}

The relationship between the derivative of the likelihood function and the
score vector is given by the following identity for each $\theta \in \Theta $%
:%
\begin{align}
s\left( X,\theta \right) & =\frac{\partial \log f\left( X,\theta \right) }{%
\partial \theta }=\frac{1}{f\left( X,\theta \right) }\frac{\partial f\left(
X,\theta \right) }{\partial \theta }  \notag \\
& \Longrightarrow \frac{\partial f\left( X,\theta \right) }{\partial \theta }%
=s\left( X,\theta \right) f\left( X,\theta \right) .  \label{2.17}
\end{align}%
We can write for each $\theta \in \Theta $: $\mathbb{E}\left[ s\left(
X,\theta \right) \right] $%
\begin{equation*}
=\int_{S}s\left( X,\theta \right) f\left( X,\theta _{0}\right) dX=\int_{S}%
\frac{1}{f\left( X,\theta \right) }\frac{\partial f\left( X,\theta \right) }{%
\partial \theta }f\left( X,\theta _{0}\right) dX.
\end{equation*}%
When the true score vector is considered, $\theta =\theta _{0}$ and the
above display yields%
\begin{equation*}
\mathbb{E}\left[ s\left( X,\theta _{0}\right) \right] =\int_{S}\frac{%
\partial f\left( X,\theta _{0}\right) }{\partial \theta }dX=\frac{\partial }{%
\partial \theta }\int_{S}f\left( X,\theta _{0}\right) dX=0
\end{equation*}

%TCIMACRO{\TeXButton{EndFrame}{\end{frame}}}%
%BeginExpansion
\end{frame}%
%EndExpansion

%TCIMACRO{\TeXButton{BeginFrame}{\begin{frame}}}%
%BeginExpansion
\begin{frame}%
%EndExpansion

\QTR{frametitle}{Score Vector and Information Matrix}

For the second part of the Theorem, use again the relationship we derived:%
\begin{gather*}
\mathbb{E}\left[ s\left( X,\theta \right) s\left( X,\theta \right) ^{\prime }%
\right] =\int_{S}s\left( X,\theta \right) s\left( X,\theta \right) ^{\prime
}f\left( X,\theta _{0}\right) dX \\
=\int_{S}\frac{\partial f\left( X,\theta \right) }{\partial \theta }s\left(
X,\theta \right) ^{\prime }\frac{f\left( X,\theta _{0}\right) }{f\left(
X,\theta \right) }dX \\
=\int_{S}\left[ \frac{\partial f\left( X,\theta \right) }{\partial \theta }%
s\left( X,\theta \right) ^{\prime }+f\left( X,\theta \right) \frac{\partial 
}{\partial \theta }s\left( X,\theta \right) ^{\prime }\right] \frac{f\left(
X,\theta _{0}\right) }{f\left( X,\theta \right) }dX \\
-\int_{S}\frac{\partial s\left( X,\theta \right) ^{\prime }}{\partial \theta 
}f\left( X,\theta _{0}\right) dX \\
=\int_{S}\frac{\partial }{\partial \theta }\left[ f\left( X,\theta \right)
s\left( X,\theta \right) ^{\prime }\right] \frac{f\left( X,\theta
_{0}\right) }{f\left( X,\theta \right) }dX-\mathbb{E}\left[ \frac{\partial }{%
\partial \theta }s\left( X,\theta \right) ^{\prime }\right]  \\
=\int_{S}\frac{\partial }{\partial \theta }\left[ f\left( X,\theta \right)
s\left( X,\theta \right) ^{\prime }\right] \frac{f\left( X,\theta
_{0}\right) }{f\left( X,\theta \right) }dX-J\left( \theta \right) 
\end{gather*}

%TCIMACRO{\TeXButton{EndFrame}{\end{frame}}}%
%BeginExpansion
\end{frame}%
%EndExpansion

%TCIMACRO{\TeXButton{BeginFrame}{\begin{frame}}}%
%BeginExpansion
\begin{frame}%
%EndExpansion

\QTR{frametitle}{Score Vector and Information Matrix}

Evaluating at $\theta =\theta _{0}$ and using (\ref{2.13}) we obtain%
\begin{eqnarray*}
\mathbb{E}\left[ s\left( X,\theta _{0}\right) s\left( X,\theta _{0}\right)
^{\prime }\right]  &=&\int_{S}\frac{\partial }{\partial \theta }\left[
f\left( X,\theta _{0}\right) s\left( X,\theta _{0}\right) ^{\prime }\right]
dX-J\left( \theta _{0}\right)  \\
&=&\frac{\partial }{\partial \theta }\int_{S}f\left( X,\theta _{0}\right)
s\left( X,\theta _{0}\right) ^{\prime }dX-J\left( \theta _{0}\right)  \\
&=&\frac{\partial }{\partial \theta }\mathbb{E}\left[ s\left( X,\theta
_{0}\right) ^{\prime }\right] -J\left( \theta _{0}\right)  \\
&=&-J\left( \theta _{0}\right) 
\end{eqnarray*}%
as required.

%TCIMACRO{\TeXButton{EndFrame}{\end{frame}}}%
%BeginExpansion
\end{frame}%
%EndExpansion

%TCIMACRO{\TeXButton{BeginFrame}{\begin{frame}}}%
%BeginExpansion
\begin{frame}%
%EndExpansion

\QTR{frametitle}{Score Vector and Information Matrix}

Note that Theorem 1 is not restricted to independent or identically
distributed samples. The above proof depends only on the validity of
Leibnitz rule for differentiation under the integral sign, ensured by (RC1)
and sufficient smoothness of the log-likelihood (RC2).

\begin{definition}
The $k\times k$ matrix%
\begin{equation}
\mathcal{I}\left( \theta \right) =-J\left( \theta \right) =-\mathbb{E}\left( 
\frac{\partial ^{2}l}{\partial \theta \partial \theta ^{\prime }}\right) 
\label{2.19}
\end{equation}%
is called (Fisher's) information matrix.
\end{definition}

%TCIMACRO{\TeXButton{EndFrame}{\end{frame}}}%
%BeginExpansion
\end{frame}%
%EndExpansion

%TCIMACRO{\TeXButton{BeginFrame}{\begin{frame}}}%
%BeginExpansion
\begin{frame}%
%EndExpansion

\QTR{frametitle}{Score Vector and Information Matrix: }

\QTR{framesubtitle}{Remarks}

\begin{enumerate}
\item Under (RC1)-(RC2), the information matrix $\mathcal{I}\left( \theta
\right) $ coincides with the covariance matrix of the score vector $s\left(
X,\theta \right) $, which ensures that $\mathcal{I}\left( \theta \right)
\geq 0$, a property of any covariance matrix. In addition, Assumption (RC3)
imposes strict negativity of the matrix $J\left( \theta \right) $ which
ensures invertability of the information matrix (see the Cram\'{e}r Rao
lower bound below).

\item In the next section we will see that, for any unbiased estimator $\hat{%
\theta}$, $var\left( \hat{\theta}\right) \geq \mathcal{I}\left( \theta
\right) ^{-1}$. If $\hat{\theta}$ is unbiased, it is centered around $\theta 
$ so the more information about $\theta $ is provided by an observation the
smaller the variance of $\hat{\theta}$ should be: hence the name
\textquotedblleft information matrix\textquotedblright\ for (\ref{2.19}).
\end{enumerate}

%TCIMACRO{\TeXButton{EndFrame}{\end{frame}}}%
%BeginExpansion
\end{frame}%
%EndExpansion

\section{The Cram\'{e}r Rao Lower Bound}

%TCIMACRO{\TeXButton{BeginFrame}{\begin{frame}}}%
%BeginExpansion
\begin{frame}%
%EndExpansion

\QTR{frametitle}{Cram\'{e}r Rao Lower Bound}

\begin{theorem}[Cram\'{e}r-Rao]
If (RC1)-(RC3) hold, then the covariance matrix of any unbiased estimator $%
\hat{\theta}$ of $\theta _{0}$\ satisfies%
\begin{equation*}
V\left( \hat{\theta}\right) \geq \mathcal{I}\left( \theta _{0}\right) ^{-1}
\end{equation*}%
in the sense that $V\left( \hat{\theta}\right) -\mathcal{I}\left( \theta
_{0}\right) ^{-1}$ is a positive semi-definite matrix.
\end{theorem}

%TCIMACRO{\TeXButton{EndFrame}{\end{frame}}}%
%BeginExpansion
\end{frame}%
%EndExpansion

%TCIMACRO{\TeXButton{BeginFrame}{\begin{frame}}}%
%BeginExpansion
\begin{frame}%
%EndExpansion

\begin{proof}
Recall that any estimator $\hat{\theta}=\hat{\theta}\left( X\right) $ is
independent of $\theta _{0}.$ Since $\hat{\theta}$ is unbiased, $\mathbb{E}%
\left( \hat{\theta}\right) =\theta _{0}$, or equivalently%
\begin{equation*}
\theta _{0}=\int_{S}\hat{\theta}\left( X\right) f\left( X,\theta _{0}\right)
dX.
\end{equation*}%
Differentiating w.r.t. $\theta _{0}^{\prime }$ we obtain%
\begin{align}
I_{k}& =\partial /\partial \theta _{0}^{\prime }\int_{S}\hat{\theta}\left(
X\right) f\left( X,\theta _{0}\right) dX  \notag \\
& =\int_{S}\hat{\theta}\left( X\right) \frac{\partial f\left( X,\theta
_{0}\right) }{\partial \theta _{0}^{\prime }}dX,\text{\quad \lbrack by (RC1)]%
}  \notag \\
& =\int_{S}\hat{\theta}\left( X\right) s\left( X,\theta _{0}\right) ^{\prime
}f\left( X,\theta _{0}\right) dX,\quad \text{[by (\ref{2.17})]}  \notag \\
& =\mathbb{E}\left[ \hat{\theta}s\left( X,\theta _{0}\right) ^{\prime }%
\right] =\mathbb{E}\left[ \left( \hat{\theta}-\theta _{0}\right) s\left(
X,\theta _{0}\right) ^{\prime }\right]  \label{2.21}
\end{align}%
since $\mathbb{E}\left[ \theta _{0}s\left( X,\theta _{0}\right) ^{\prime }%
\right] =\theta _{0}\mathbb{E}\left[ s\left( X,\theta _{0}\right) ^{\prime }%
\right] =0$ by (\ref{2.14}).
\end{proof}

%TCIMACRO{\TeXButton{EndFrame}{\end{frame}}}%
%BeginExpansion
\end{frame}%
%EndExpansion

%TCIMACRO{\TeXButton{BeginFrame}{\begin{frame}}}%
%BeginExpansion
\begin{frame}%
%EndExpansion

\QTR{frametitle}{Cram\'{e}r Rao Lower Bound}

\begin{proof}
First, assume that $k=1$ ($\theta _{0}$ is a scalar) for simplicity. Then
both $s\left( X,\theta _{0}\right) $ and $\mathcal{I}\left( \theta
_{0}\right) $ are scalars and the equality above reduces to%
\begin{align*}
1& =\mathbb{E}\left[ \left( \hat{\theta}-\theta _{0}\right) s\left( X,\theta
_{0}\right) \right] =cov\left( \hat{\theta},s\left( X,\theta _{0}\right)
\right) \\
& \leq \left[ var\left( \hat{\theta}\right) \right] ^{1/2}\left[ var\left(
s\left( X,\theta _{0}\right) \right) \right] ^{1/2} \\
& =\left[ var\left( \hat{\theta}\right) \right] ^{1/2}\left[ \mathcal{I}%
\left( \theta _{0}\right) \right] ^{1/2}.
\end{align*}%
by Cauchy-Schwarz inequality.

This shows the result for scalar $\theta _{0}$.

The key property that delivers the lower bound is the CS inequality.

Note that in this case the ordering is complete, i.e. $var\left( \hat{\theta}%
\right) \geq \mathcal{I}\left( \theta _{0}\right) ^{-1}$ holds as a true
inequality between numbers.
\end{proof}

%TCIMACRO{\TeXButton{EndFrame}{\end{frame}}}%
%BeginExpansion
\end{frame}%
%EndExpansion

%TCIMACRO{\TeXButton{BeginFrame}{\begin{frame}}}%
%BeginExpansion
\begin{frame}%
%EndExpansion

\begin{proof}
For the general case, the standard Cauchy-Schwarz inequality is not
sufficient; a matrix generalisation of the form%
\begin{equation}
\mathbb{E}\left( xy^{\prime }\right) \left[ \mathbb{E}\left( yy^{\prime
}\right) \right] ^{-1}\mathbb{E}\left( yx^{\prime }\right) \leq \mathbb{E}%
\left( xx^{\prime }\right)  \label{mCS}
\end{equation}%
(in the positive semi-definite sense) for all $k\times 1$ vectors $x,y$ with
finite second moment and $\mathbb{E}\left( yy^{\prime }\right) >0$.

Letting $x=\hat{\theta}-\theta _{0}$ and $y=s\left( X,\theta _{0}\right) $, (%
\ref{mCS}) yields%
\begin{gather*}
\underset{=I_{k}}{\underbrace{\mathbb{E}\left\{ \left( \hat{\theta}-\theta
_{0}\right) s\left( X,\theta _{0}\right) ^{\prime }\right\} }}[\underset{=%
\mathcal{I}\left( \theta _{0}\right) }{\underbrace{\mathbb{E}\left\{ s\left(
X,\theta _{0}\right) s\left( X,\theta _{0}\right) ^{\prime }\right\} }}]^{-1}
\\
\underset{=I_{k}}{\times \underbrace{\mathbb{E}\left[ s\left( X,\theta
_{0}\right) \left( \hat{\theta}-\theta _{0}\right) ^{\prime }\right] }}\leq 
\underset{=V\left( \hat{\theta}\right) }{\underbrace{\mathbb{E}\left[ \left( 
\hat{\theta}-\theta _{0}\right) \left( \hat{\theta}-\theta _{0}\right)
^{\prime }\right] }}
\end{gather*}%
we obtain $\mathcal{I}\left( \theta _{0}\right) ^{-1}\leq V\left( \hat{\theta%
}\right) $.
\end{proof}

%TCIMACRO{\TeXButton{EndFrame}{\end{frame}}}%
%BeginExpansion
\end{frame}%
%EndExpansion

%TCIMACRO{\TeXButton{BeginFrame}{\begin{frame}}}%
%BeginExpansion
\begin{frame}%
%EndExpansion

\QTR{frametitle}{Cram\'{e}r Rao Lower Bound}

Theorem 2 implies that the MSE of any unbiased estimator $\hat{\theta}$ of $%
\theta _{0}$ cannot be \textquotedblleft smaller\textquotedblright\ than $%
\mathcal{I}\left( \theta _{0}\right) ^{-1}$, i.e. the best possible unbiased
estimator is one whose covariance matrix achieves the Cram\'{e}r Rao lower
bound.

Such estimators are called efficient.

\begin{definition}
An unbiased estimator $\hat{\theta}$ of $\theta _{0}$ is called \textbf{%
efficient} if 
\begin{equation*}
V\left( \hat{\theta}\right) =\mathcal{I}\left( \theta _{0}\right) ^{-1}.
\end{equation*}
\end{definition}

%TCIMACRO{\TeXButton{EndFrame}{\end{frame}}}%
%BeginExpansion
\end{frame}%
%EndExpansion

\section{Consistency of the MLE}

%TCIMACRO{\TeXButton{BeginFrame}{\begin{frame}}}%
%BeginExpansion
\begin{frame}%
%EndExpansion

\QTR{frametitle}{Consistency of the MLE}

\begin{itemize}
\item The main problem with deriving the asymptotic properties of the MLE
estimator is that, unlike linear estimation procedures such as OLS, there is
no general closed form expression for the MLE.

\item This is always the case in non-linear estimation procedures based on
maximising an objective function since, even when the first order conditions
are simple enough to deliver an estimator in closed form, this form will
vary according to the objective function (e.g. different likelihood
functions give rise to different MLEs).

\item Ultimately, limit distribution theory will be derived by a
linearisation scheme based on the mean value theorem (see the next section),
but successful application of this methods requires first to establish the
consistency of the MLE.
\end{itemize}

%TCIMACRO{\TeXButton{EndFrame}{\end{frame}}}%
%BeginExpansion
\end{frame}%
%EndExpansion

%TCIMACRO{\TeXButton{BeginFrame}{\begin{frame}}}%
%BeginExpansion
\begin{frame}%
%EndExpansion

\QTR{frametitle}{Consistency of the MLE}

\QTR{framesubtitle}{Assumptions}

\begin{description}
\item[(C1)] The sample random vector $X=\left( X_{1},...,X_{n}\right) $
consists of \textbf{i.i.d.} random variables with marginal density $f\left(
X_{1},\theta \right) $ that is continuous w.r.t. $\theta $.

\item[(C2)] $\Theta $ is a compact (equivalently closed and bounded) subset
of $\mathbb{R}^{k}$.

\item[(C3)] $\theta _{0}$ lies on the interior of $\Theta $ and is the 
\textit{unique} maximiser\footnote{%
We know that $\theta _{0}$ is a maximiser of $E\log f\left( X_{1},\theta
\right) $.} of $\mathbb{E}\log f\left( X_{1},\theta \right) $ over $\Theta $%
. In other words,%
\begin{equation}
\kappa \left( \theta \right) :=\mathbb{E}\log f\left( X_{1},\theta
_{0}\right) -\mathbb{E}\log f\left( X_{1},\theta \right) >0  \label{kappa}
\end{equation}%
for all $\theta \in \Theta \backprime \left\{ \theta _{0}\right\} $.

\item[(C4)] The i.i.d. sequence $\left\{ \log f\left( X_{i},\theta \right)
:i\geq 1\right\} $ satisfies a \textit{uniform} law of large numbers over $%
\Theta $:%
\begin{equation}
\max_{\theta \in \Theta }\left\vert \frac{1}{n}\sum_{i=1}^{n}\log f\left(
X_{i},\theta \right) -\mathbb{E}\left[ \log f\left( X_{1},\theta \right) %
\right] \right\vert \rightarrow _{p}0\;\text{{\small as} }n\rightarrow
\infty .  \label{ULLNc}
\end{equation}
\end{description}

%TCIMACRO{\TeXButton{EndFrame}{\end{frame}}}%
%BeginExpansion
\end{frame}%
%EndExpansion

%TCIMACRO{\TeXButton{BeginFrame}{\begin{frame}}}%
%BeginExpansion
\begin{frame}%
%EndExpansion

\QTR{frametitle}{Consistency of the MLE}

\begin{itemize}
\item Compactness of the parameter space $\Theta $ together with the
continuity of the log-likelihood function%
\begin{equation}
\log f\left( X,\theta \right) =\sum_{i=1}^{n}\log f\left( X_{i},\theta
\right)  \label{lL}
\end{equation}%
(ensured by C1) imply that the MLE $\hat{\theta}_{n}$ belongs to $\Theta $
for each $n$ (by Weierstrass Theorem).

\item For the same reason (compactness of $\Theta $ and continuity of $%
f\left( X_{i},\cdot \right) $) we can write $\max $ instead of $\sup $ in (%
\ref{ULLNc}).
\end{itemize}

%TCIMACRO{\TeXButton{EndFrame}{\end{frame}}}%
%BeginExpansion
\end{frame}%
%EndExpansion

%TCIMACRO{\TeXButton{BeginFrame}{\begin{frame}}}%
%BeginExpansion
\begin{frame}%
%EndExpansion

\QTR{frametitle}{Consistency of the MLE}

\begin{itemize}
\item For (C4), note that 
\begin{equation*}
n^{-1}\sum_{i=1}^{n}\log f\left( X_{i},\theta \right) \rightarrow _{p}%
\mathbb{E}\left[ \log f\left( X_{1},\theta \right) \right]
\end{equation*}%
as $n\rightarrow \infty $ for each $\theta $ by a standard WLLN.

\item However, we require a uniform version of this WLLN over $\Theta .$ The
validity of (\ref{ULLNc}) can be reduced to more elementary conditions
(uniform stochastic equicontinuity) but the technical requirements are far
beyond the scope of this course.
\end{itemize}

%TCIMACRO{\TeXButton{EndFrame}{\end{frame}}}%
%BeginExpansion
\end{frame}%
%EndExpansion

%TCIMACRO{\TeXButton{BeginFrame}{\begin{frame}}}%
%BeginExpansion
\begin{frame}%
%EndExpansion

\QTR{frametitle}{Consistency of the MLE}

\begin{theorem}[Consistency]
Under (C1)-(C4), $\hat{\theta}_{n}\rightarrow _{p}\theta _{0}$ as $%
n\rightarrow \infty .$
\end{theorem}

We will not prove the result.

%TCIMACRO{\TeXButton{EndFrame}{\end{frame}}}%
%BeginExpansion
\end{frame}%
%EndExpansion

\section{Asymptotic Normality of the MLE}

%TCIMACRO{\TeXButton{BeginFrame}{\begin{frame}}}%
%BeginExpansion
\begin{frame}%
%EndExpansion

\QTR{frametitle}{Asymptotic Normality of the MLE}

Denote by $f_{n}\left( X,\theta \right) $ the likelihood function of a
sample $X=\left( X_{1},...,X_{n}\right) $ with $\theta _{0}$ being the true
value of $\theta $, by $\hat{\theta}_{n}$ the MLE of $\theta _{0}$, by $%
s_{n}\left( \theta \right) $ the score vector and by $\mathcal{I}_{n}\left(
\theta \right) $ the information matrix, and finally let%
\begin{equation}
z_{i}\left( \theta \right) :=\frac{\partial \log f\left( X_{i},\theta
\right) }{\partial \theta }\;\;\;\text{and\ }\;\;H_{i}\left( \theta \right)
:=\frac{\partial ^{2}\log f\left( X_{i},\theta \right) }{\partial \theta
\partial \theta ^{\prime }}  \label{zH}
\end{equation}%
denote the "marginal" score and Hessian.

If the sample vector $X=\left( X_{1},...,X_{n}\right) $ consists of i.i.d.
random variables, then $\left\{ H_{i}\left( \theta \right)
:i=1,...,n\right\} $ are i.i.d. random matrices and%
\begin{align}
\mathcal{I}_{n}\left( \theta \right) & =-\mathbb{E}\left( \frac{\partial
^{2}\log f_{n}\left( X,\theta \right) }{\partial \theta \partial \theta
^{\prime }}\right) =-\mathbb{E}\left[ \frac{\partial ^{2}}{\partial \theta
\partial \theta ^{\prime }}\sum_{i=1}^{n}\log f\left( X_{i},\theta \right) %
\right]   \notag \\
& =-\sum_{i=1}^{n}\mathbb{E}\left( \frac{\partial ^{2}\log f\left(
X_{i},\theta \right) }{\partial \theta \partial \theta ^{\prime }}\right) =-n%
\mathbb{E}\left[ H_{1}\left( \theta \right) \right] .  \label{2.20}
\end{align}

%TCIMACRO{\TeXButton{EndFrame}{\end{frame}}}%
%BeginExpansion
\end{frame}%
%EndExpansion

%TCIMACRO{\TeXButton{BeginFrame}{\begin{frame}}}%
%BeginExpansion
\begin{frame}%
%EndExpansion

\QTR{frametitle}{Asymptotic Normality of the MLE}

\begin{itemize}
\item Note that $z_{i}\left( \theta \right) $ and $H_{i}\left( \theta
\right) $ in (\ref{zH}) are\textit{\ independent of} $n$ (why?) and, hence,
so is the matrix $\mathbb{E}\left[ H_{1}\left( \theta \right) \right] $ in (%
\ref{2.20}). The identity (\ref{2.20}) implies that the information matrix $%
\mathcal{I}_{n}\left( \theta \right) $ increases to $\infty $ with rate
equal to the sample size $n$ as $n\rightarrow \infty $:%
\begin{equation}
\lim_{n\rightarrow \infty }\frac{1}{n}\mathcal{I}_{n}\left( \theta \right) =-%
\mathbb{E}\left[ H_{1}\left( \theta \right) \right] .  \label{limI}
\end{equation}

\item While our derivation of (\ref{2.20}) is confined to the case where $%
X_{1},...,X_{n}$ are i.i.d. random variables, the limit behaviour of the
information matrix is given by (\ref{limI}) in much more general situations:
it applies even under dependent sampling schemes.
\end{itemize}

%TCIMACRO{\TeXButton{EndFrame}{\end{frame}}}%
%BeginExpansion
\end{frame}%
%EndExpansion

%TCIMACRO{\TeXButton{BeginFrame}{\begin{frame}}}%
%BeginExpansion
\begin{frame}%
%EndExpansion

\QTR{frametitle}{Asymptotic Normality of the MLE}

\begin{itemize}
\item A standard way to address this non-linearity problem is to linearise
the derivative of the objective function (in our case the score) around the
estimator using a mean value theorem between the estimator $\hat{\theta}_{n}$
and the true value $\theta _{0}$.

\item Since the score at $\hat{\theta}_{n}$ is zero, one of the constant
terms in the linearised form disappears and we are left with a linear
relationship between the score at $\theta _{0}$ and $\hat{\theta}_{n}-\theta
_{0}$. Here is this method in action:

\item \textbf{Mean value theorem on }$s_{n}\left( \theta _{0}\right) $%
\textbf{\ around }$\hat{\theta}_{n}$\textbf{:}%
\begin{equation}
s_{n}\left( \theta _{0}\right) =s_{n}\left( \hat{\theta}_{n}\right) +\left. 
\frac{\partial }{\partial \theta }s_{n}\left( \theta \right) ^{\prime
}\right\vert _{\theta =\theta _{n}^{\ast }}\left( \theta _{0}-\hat{\theta}%
_{n}\right)   \label{MVT}
\end{equation}%
where $\theta _{n}^{\ast }$ is intermediate to $\hat{\theta}_{n}$ and $%
\theta _{0}$:%
\begin{equation}
\left\Vert \theta _{n}^{\ast }-\theta _{0}\right\Vert \leq \left\Vert \hat{%
\theta}_{n}-\theta _{0}\right\Vert .  \label{interm}
\end{equation}
\end{itemize}

%TCIMACRO{\TeXButton{EndFrame}{\end{frame}}}%
%BeginExpansion
\end{frame}%
%EndExpansion

%TCIMACRO{\TeXButton{BeginFrame}{\begin{frame}}}%
%BeginExpansion
\begin{frame}%
%EndExpansion

\QTR{frametitle}{Asymptotic Normality of the MLE}

Since $s_{n}\left( \hat{\theta}_{n}\right) =0$ by the first order condition
and%
\begin{eqnarray*}
\left. \frac{\partial }{\partial \theta }s_{n}\left( \theta \right) ^{\prime
}\right\vert _{\theta =\theta _{n}^{\ast }} &=&\left. \frac{\partial }{%
\partial \theta }\frac{\partial \log f_{n}\left( X,\theta \right) }{\partial
\theta ^{\prime }}\right\vert _{\theta =\theta _{n}^{\ast }}= \\
\frac{\partial ^{2}\log f_{n}\left( X,\theta _{n}^{\ast }\right) }{\partial
\theta \partial \theta ^{\prime }} &=&\sum_{i=1}^{n}\frac{\partial ^{2}\log
f\left( X_{i},\theta _{n}^{\ast }\right) }{\partial \theta \partial \theta
^{\prime }} \\
&=&\sum_{i=1}^{n}H_{i}\left( \theta _{n}^{\ast }\right) 
\end{eqnarray*}%
in the notation of (\ref{zH}), we can write (\ref{MVT}) as%
\begin{equation}
s_{n}\left( \theta _{0}\right) =-\sum_{i=1}^{n}H_{i}\left( \theta _{n}^{\ast
}\right) \left( \hat{\theta}_{n}-\theta _{0}\right) .  \label{MVT1}
\end{equation}

%TCIMACRO{\TeXButton{EndFrame}{\end{frame}}}%
%BeginExpansion
\end{frame}%
%EndExpansion

%TCIMACRO{\TeXButton{BeginFrame}{\begin{frame}}}%
%BeginExpansion
\begin{frame}%
%EndExpansion

\QTR{frametitle}{Asymptotic Normality of the MLE}

The promised linearisation is established by (\ref{MVT1}). The last step is
to write the score in terms of a sum of random vectors:%
\begin{equation}
s_{n}\left( \theta _{0}\right) =\frac{\partial \log f_{n}\left( X,\theta
_{0}\right) }{\partial \theta }=\sum_{i=1}^{n}\frac{\partial \log f\left(
X_{i},\theta _{0}\right) }{\partial \theta }=\sum_{i=1}^{n}z_{i}\left(
\theta _{0}\right)   \label{Z}
\end{equation}%
in the notation of (\ref{zH}). Substituting into (\ref{MVT1}), we obtain%
\begin{equation}
\sum_{i=1}^{n}z_{i}\left( \theta _{0}\right) =-\sum_{i=1}^{n}H_{i}\left(
\theta _{n}^{\ast }\right) \left( \hat{\theta}_{n}-\theta _{0}\right) .
\label{MVT2}
\end{equation}

%TCIMACRO{\TeXButton{EndFrame}{\end{frame}}}%
%BeginExpansion
\end{frame}%
%EndExpansion

%TCIMACRO{\TeXButton{BeginFrame}{\begin{frame}}}%
%BeginExpansion
\begin{frame}%
%EndExpansion

\QTR{frametitle}{Asymptotic Normality Conditions}

\begin{description}
\item[(N1)] The sample random vector $X=\left( X_{1},...,X_{n}\right) $
consists of \textbf{i.i.d.} r.v.s with marginal density $f\left(
X_{1},\theta \right) $ that is twice continuously differentiable w.r.t. $%
\theta $.

\item[(N2)] The support of $f\left( X_{1},\theta \right) $, $\left\{ X\in 
\mathbb{R}:f\left( X_{1},\theta \right) >0\right\} $, is independent of $%
\theta $.

\item[(N3)] $\mathbb{E}\left\Vert H_{1}\left( \theta \right) \right\Vert
<\infty $ for all $\theta \in \Theta $\ and $\mathbb{E}\left[ H_{1}\left(
\theta _{0}\right) \right] $ is negative definite.

\item[(N4)] The MLE $\hat{\theta}_{n}$ is consistent.

\item[(N5)] The sequence of random matrices $H_{i}\left( \theta \right) $ in
(\ref{zH}) satisfies the following \textit{uniform} law of large numbers on
an open neighbourhood of $\theta _{0}$ $N\left( \delta \right) =\left\{
\theta \in \Theta :\left\Vert \theta -\theta _{0}\right\Vert <\delta
\right\} $: there exists $\delta >0$ such that%
\begin{equation}
\sup_{\theta \in N\left( \delta \right) }\left\Vert \frac{1}{n}%
\sum_{i=1}^{n}H_{i}\left( \theta \right) -\mathbb{E}\left[ H_{1}\left(
\theta _{0}\right) \right] \right\Vert \rightarrow _{p}0\;\;\;\;\text{as }%
n\rightarrow \infty .  \label{ULLN}
\end{equation}
\end{description}

%TCIMACRO{\TeXButton{EndFrame}{\end{frame}}}%
%BeginExpansion
\end{frame}%
%EndExpansion

%TCIMACRO{\TeXButton{BeginFrame}{\begin{frame}}}%
%BeginExpansion
\begin{frame}%
%EndExpansion

\QTR{frametitle}{Asymptotic Normality of the MLE}

\begin{itemize}
\item Conditions (N2) and (N3) is a restatement of (RC2) and (RC3).

\item Uniformity in (N5) serves the purpose of dealing with the intermediate
point $\theta _{n}^{\ast }$ arising from the linearisation scheme: $%
n^{-1}\sum_{i=1}^{n}H_{i}\left( \theta \right) \rightarrow _{p}\mathbb{E}%
\left[ H_{1}\left( \theta \right) \right] $ as $n\rightarrow \infty $ for
each $\theta $ by a standard WLLN since $\left\{ H_{i}\left( \theta \right)
\right\} $ is an i.i.d. sequence with finite first moment.

\item However, this WLLN is not enough as, looking at (\ref{MVT2}), it is
clear that a law of large numbers on $n^{-1}\sum_{i=1}^{n}H_{i}\left( \theta
_{n}^{\ast }\right) $ is required.
\end{itemize}

%TCIMACRO{\TeXButton{EndFrame}{\end{frame}}}%
%BeginExpansion
\end{frame}%
%EndExpansion

%TCIMACRO{\TeXButton{BeginFrame}{\begin{frame}}}%
%BeginExpansion
\begin{frame}%
%EndExpansion

\QTR{frametitle}{Asymptotic Normality of the MLE}

\begin{lemma}
Under (N1)-(N3), the random vectors $\left\{ z_{i}\left( \theta _{0}\right)
:1\leq i\leq n\right\} $ in (\ref{zH}) are i.i.d. with $\mathbb{E}%
z_{1}\left( \theta _{0}\right) =0$ and $\mathbb{E}\left[ z_{1}\left( \theta
_{0}\right) z_{1}\left( \theta _{0}\right) ^{\prime }\right] =-\mathbb{E}%
\left[ H_{1}\left( \theta \right) \right] $. In particular, as $n\rightarrow
\infty $:%
\begin{equation}
\frac{1}{n}\sum_{i=1}^{n}z_{i}\left( \theta _{0}\right) \rightarrow
_{p}0\;\;\;\text{and}\;\;\;\frac{1}{\sqrt{n}}\sum_{i=1}^{n}z_{i}\left(
\theta _{0}\right) \rightarrow _{d}N\left( 0,-\mathbb{E}\left[ H_{1}\left(
\theta _{0}\right) \right] \right) .  \label{limZ}
\end{equation}
\end{lemma}

%TCIMACRO{\TeXButton{EndFrame}{\end{frame}}}%
%BeginExpansion
\end{frame}%
%EndExpansion

%TCIMACRO{\TeXButton{BeginFrame}{\begin{frame}}}%
%BeginExpansion
\begin{frame}%
%EndExpansion

\begin{proof}
By their definition in (\ref{Z}), $\left\{ z_{i}\left( \theta _{0}\right)
:1\leq i\leq n\right\} $ are i.i.d. because the sample $X=\left(
X_{1},...,X_{n}\right) $ consists of i.i.d. random variables $X_{i}$. We
evaluate the first two moments of $z_{1}\left( \theta _{0}\right) $: by (\ref%
{Z}), the theorem on the first two moments of the score vector and the
identical distribution assumption on $X_{i}$,%
\begin{equation*}
0=\mathbb{E}s_{n}\left( \theta _{0}\right) =\sum\nolimits_{i=1}^{n}\mathbb{E}%
z_{i}\left( \theta _{0}\right) =n\mathbb{E}z_{1}\left( \theta _{0}\right)
\Longrightarrow \mathbb{E}z_{1}\left( \theta _{0}\right) =0.
\end{equation*}%
Similarly, the theorem on the score vector and the i.i.d. property of $%
z_{i}\left( \theta _{0}\right) $ give%
\begin{gather*}
-n\mathbb{E}\left[ H_{1}\left( \theta _{0}\right) \right] =\mathcal{I}%
_{n}\left( \theta _{0}\right) =\mathbb{E[}s_{n}\left( \theta _{0}\right)
s_{n}\left( \theta _{0}\right) ^{\prime }] \\
=\sum\nolimits_{i=1}^{n}\sum\nolimits_{j=1}^{n}\mathbb{E}\left[ z_{i}\left(
\theta _{0}\right) z_{j}\left( \theta _{0}\right) ^{\prime }\right] \\
=\sum\nolimits_{i=1}^{n}\mathbb{E}\left[ z_{i}\left( \theta _{0}\right)
z_{i}\left( \theta _{0}\right) ^{\prime }\right] =n\mathbb{E}\left[
z_{1}\left( \theta _{0}\right) z_{1}\left( \theta _{0}\right) ^{\prime }%
\right]
\end{gather*}
\end{proof}

%TCIMACRO{\TeXButton{EndFrame}{\end{frame}}}%
%BeginExpansion
\end{frame}%
%EndExpansion

%TCIMACRO{\TeXButton{BeginFrame}{\begin{frame}}}%
%BeginExpansion
\begin{frame}%
%EndExpansion

\begin{proof}
So, since $\mathcal{I}_{n}\left( \theta _{0}\right) =-n\mathbb{E}\left[
H_{1}\left( \theta _{0}\right) \right] $ by (\ref{2.20}), we obtain $\mathbb{%
E}\left[ z_{1}\left( \theta _{0}\right) z_{1}\left( \theta _{0}\right)
^{\prime }\right] =-\mathbb{E}\left[ H_{1}\left( \theta _{0}\right) \right] $%
.

Since $\left\{ z_{i}\left( \theta _{0}\right) :1\leq i\leq n\right\} $ are
i.i.d. random vectors, the first part of (\ref{limZ}) follows by the
standard WLLN and the second part of (\ref{limZ}) follows by the standard
CLT.
\end{proof}

%TCIMACRO{\TeXButton{EndFrame}{\end{frame}}}%
%BeginExpansion
\end{frame}%
%EndExpansion

%TCIMACRO{\TeXButton{BeginFrame}{\begin{frame}}}%
%BeginExpansion
\begin{frame}%
%EndExpansion

\QTR{frametitle}{Asymptotic Normality of the MLE}

\begin{theorem}
Under (N1)-(N5),%
\begin{equation*}
\sqrt{n}\left( \hat{\theta}_{n}-\theta _{0}\right) \rightarrow _{d}\mathcal{N%
}\left( 0,-\mathbb{E}\left[ H_{1}\left( \theta _{0}\right) \right]
^{-1}\right) .
\end{equation*}
\end{theorem}

%TCIMACRO{\TeXButton{EndFrame}{\end{frame}}}%
%BeginExpansion
\end{frame}%
%EndExpansion

%TCIMACRO{\TeXButton{BeginFrame}{\begin{frame}}}%
%BeginExpansion
\begin{frame}%
%EndExpansion

\QTR{frametitle}{Asymptotic Normality of the MLE}

\begin{proof}
Since $\mathbb{E}\left[ H_{1}\left( \theta _{0}\right) \right] <0$ by (C3),
we conclude that the matrix $\frac{1}{n}\sum_{i=1}^{n}H_{i}\left( \theta
_{n}^{\ast }\right) $ is negative definite for $n$ large enough. Rearranging
(\ref{MVT2}) then yields for $n$ large enough%
\begin{eqnarray*}
\sqrt{n}\left( \hat{\theta}_{n}-\theta _{0}\right) &=&\left[ -\frac{1}{n}%
\sum_{i=1}^{n}H_{i}\left( \theta _{n}^{\ast }\right) \right] ^{-1}\frac{1}{%
\sqrt{n}}\sum_{i=1}^{n}z_{i}\left( \theta _{0}\right) \\
&\rightarrow &\!\!\!\!\!_{d}\text{ \ }-\mathbb{E}\left[ H_{1}\left( \theta
_{0}\right) \right] ^{-1}\mathcal{N}\left( 0,-\mathbb{E}\left[ H_{1}\left(
\theta _{0}\right) \right] \right) \\
&\rightarrow &\!\!\!\!\!_{d}\text{ \ }\mathcal{N}\left( 0,-\mathbb{E}\left[
H_{1}\left( \theta _{0}\right) \right] ^{-1}\right)
\end{eqnarray*}%
as required.
\end{proof}

%TCIMACRO{\TeXButton{EndFrame}{\end{frame}}}%
%BeginExpansion
\end{frame}%
%EndExpansion

\section{Asymptotic Efficiency}

%TCIMACRO{\TeXButton{BeginFrame}{\begin{frame}}}%
%BeginExpansion
\begin{frame}%
%EndExpansion

\QTR{frametitle}{Asymptotic Efficiency}

\begin{corollary}
The MLE $\hat{\theta}_{n}$ attains the Cram\'{e}r-Rao lower bound as $%
n\rightarrow \infty $.
\end{corollary}

\begin{proof}
Theorem 4 implies that%
\begin{equation*}
V\left( \hat{\theta}_{n}\right) =\frac{1}{n}V\left[ \sqrt{n}\left( \hat{%
\theta}_{n}-\theta _{0}\right) \right] \sim -\frac{1}{n}\mathbb{E}\left[
H_{1}\left( \theta _{0}\right) \right] ^{-1}
\end{equation*}%
as $n\rightarrow \infty $. By (\ref{2.20}), $\mathcal{I}_{n}\left( \theta
_{0}\right) =-n\mathbb{E}\left[ H_{1}\left( \theta _{0}\right) \right] $,
giving%
\begin{equation*}
\lim_{n\rightarrow \infty }\left[ V\left( \hat{\theta}_{n}\right) -\mathcal{I%
}_{n}\left( \theta _{0}\right) ^{-1}\right] =0.
\end{equation*}%
We conclude that: \textit{since the MLE is asymptotically unbiased}, it is
the \textbf{best} possible estimator in large samples (in MSE sense).
\end{proof}

%TCIMACRO{\TeXButton{EndFrame}{\end{frame}}}%
%BeginExpansion
\end{frame}%
%EndExpansion

%TCIMACRO{\TeXButton{BeginFrame}{\begin{frame}}}%
%BeginExpansion
\begin{frame}%
%EndExpansion

\QTR{frametitle}{MLE Applications}

\begin{itemize}
\item There are various MLE applications

\item In economics: binary dependent variable models: Logit, Probit, Tobit...

\item In finance: ARCH, GARCH...

\item Stochastic volatility models, state space models

\item Bayesian econometrics
\end{itemize}

%TCIMACRO{\TeXButton{EndFrame}{\end{frame}}}%
%BeginExpansion
\end{frame}%
%EndExpansion

%TCIMACRO{\TeXButton{BeginFrame}{\begin{frame}}}%
%BeginExpansion
\begin{frame}%
%EndExpansion

\QTR{frametitle}{Quasi MLE}

\begin{itemize}
\item The most crucial assumption that we did not explicitly laid out is
that the density of the data $f(X;\theta )$ is known.

\item MLE\ is a full-information technique and is more efficient than
limited-information methods only under correctly specified $f(X;\theta ).$

\item It turns out, however, that if $f(X;\theta )$ is mis-specified,
consistency and asymptotic normality can be obtained nevertheless.

\item Such methods are called Quasi- or Pseudo- MLE.

\item In particular, there exists a large number of pseudo true
distributions leading to consistent QMLEs.

\item The cost that one needs to bear is loss of efficiency,%
\begin{equation*}
\sqrt{n}\left( \hat{\theta}_{n}-\theta _{0}\right) \rightarrow \!_{d}\text{
\ }\mathcal{N}\left( 0,I(\theta _{0})^{-1}\mathbb{E}_{Y}\mathbb{E}%
_{0}[s(X,\theta _{0})s(X,\theta _{0})^{\prime }]I(\theta _{0})^{-1}\right)
\end{equation*}

where $\mathbb{E}_{Y}\mathbb{E}_{0}[s(X,\theta _{0})s(X,\theta _{0}^{\prime
})]\neq I(\theta _{0})=-\mathbb{E}_{Y}\mathbb{E}_{0}J(X,\theta _{0}).$
\end{itemize}

%TCIMACRO{\TeXButton{EndFrame}{\end{frame}}}%
%BeginExpansion
\end{frame}%
%EndExpansion

\end{document}
